\chapter{Introduction} 

To move an aircraft, a propulsion system is needed to create thrust. Different engines develop thrust in different ways, but all depend on burning fuels. The chemical process of burning a fuel is called combustion. In such a process, oxygen reacts rapidly with a fuel which can be solid, liquid, or gas. During combustion, new chemical substances, called exhaust, are produced, most of which are water (\ce{H2O}, hydrogen + oxygen) and carbon dioxide (\ce{CO2}, carbon + oxygen). However, the emission of by-products are inevitable such as nitrogen oxides (\ce{NO_x}, nitrogen + oxygen) since nitrogen constitutes around 78\% of air besides the presence of dioxygen. Nitric oxide (\ce{NO}) and nitrogen dioxide (\ce{NO2}) contribute to the formation of smog and acid rain and badly influences ozone. Another air pollutant gas, carbon monoxide (\ce{CO}), also plays roles affecting climate change. Therefore, adequate tools are required for estimating emissions. This thesis aims at developing a simulator written in the Julia programming language following the chemical reactor network approach to model the formation of pollutants. In this project, we focus on the combustion of natural gas with methane (\ce{CH4}) as its main component\footnote{we neglect here the presence of higher alkanes such as propane (\ce{C3H8}) and butane (\ce{C4H10}).} in large industrial furnaces. 

An overview is given in this chapter. Chapter~\ref{ch:combustion} further describes the combustion of methane and its reaction network in furnaces. Chapter~\ref{ch:julia} introduces the ecosystem of the Julia programming language. {\large \color{red} TODO}